\label{ch:tp-bus}
\chapter{Sběrnice a komunikace}

Jak už bylo zmíněno v kapitole o CPU, sběrnice je hlavní prostředek pro zajištění fyzické a logické
komunikaci mezi procesorem a jinými zařízeními.

V závislosti od složitosti návrhu architektury a možností počítače, může architektura obsahovat i
více sběrnic. Například, jak již bylo zmíněno v sekci o architektonických modelech, v Harvardovým
modelu instrukce a data nejsou uchovávány ve stejné pamětí a mají rozdílné, nezávislé sběrnice. Také
moderní stolní počítače disponují uživatelsky dostupnými sběrnicemi jako např. USB, PCIe, FireWire,
Ethernet apod., které slouží k fyzickému přípojení externích vstupných a výstupních zařízení, nebo
jiných spotřebních zařízení.

Pokud počítač obsahuje jednu sběrnici, která propojuje většinu majoritních komponentů, tak se nazývá
\textbf{systémová sběrnice}, a je rozdělena do 3 funkčních skupin signálů (menších sběrnic):

\begin{itemize}
    \item \textbf{Adresová sběrnice} \textit{(address bus)} zadává paměťovou
        adresu (každý její vedení odpovídá bitu paměťové adresy).
    \item \textbf{Datová sběrnice} \textit{(data bus)} je určená k přenosu dat mezi procesorem a
        destinaci, jejíž adresa byla nastavena na adresové sběrnici.
    \item \textbf{Řídicí sběrnice} \textit{(control bus)} určuje průběh komunikaci, tj. jestli se
        jedná o operaci čtení nebo zápisu, časování přenosu atd.
\end{itemize}

Každá sběrnice má určitou tzv. \textbf{šířku}, která určuje velikost dat kterou je schopná přenést.
V případě adresové sběrnice, šířka přímo ovlivňuje velikost adresovatelné pamětí procesorem.

\section{Princip komunikace}

Princip komunikace spočívá v tom, že veškerou komunikaci zahajuje vždy jenom jedno zařízení (typicky
CPU), běžně označované jako \textbf{master}. Jakékoli zařízení, se kterým master zahájí komunikaci,
je označováno jako \textbf{slave} (např. paměť, I/O zařízení apod.). Master nastaví adresu zařízení
na adresové sběrnici, a také odpovídající konfiguraci řídicí sběrnice (zda se jedna o čtení nebo
zápis apod.), a slave odpovídá na požadavky mastera. Slave sleduje adresovou sběrnici a pokud
rozpozná svou adresu, odpoví (pošle data nebo odpoví).

U složitějších systému na sběrnici může být i více masterů najednou, a tento mechanizmus se nazývá
\textbf{arbitráž sběrnice}, aby předejit konfliktům.

Typicky pro komunikaci se používají mapované registry (viz. další sekce). Zařízení mapuje do
adresního prostoru CPU registry, které dovolují procesoru komunikovat a ovládat mapované zařízení.

\section{Dekódování adres a mapování}

Jak již bylo zmíněno v kapitole o pamětí, procesor (a programátor) vidí paměť jako jednolitý
adresový prostor. Fyzicky tento prostor může být rozdělen na oblastí, které jsou přiřazená externím
zařízením. Proces výběru odpovídajícího zařízení na základě adresy se nazývá \textbf{dekódování
adresy}.

\begin{figure}[H]
    \centering
    \includesvg[width=\textwidth]{images/memory_map_ex.svg}
    \caption{Podrobní znázornění způsobu mapování periferií (neexistující systém).}
    \label{fig:memory_map_ex}
\end{figure}

V starších systémech kvůli šetření na logice dekódování, častým jevem bylo \textbf{neúplné dekódování
adres}, kdy logika dekódování porovnávala jenom některé bity adresy, což vedlo k tomu, že to a samé
zařízení mohlo být přistupováno pomocí více adres. Tento jev se nazývá \textbf{zrcadlení}
(\textit{mirroring}).

\textbf{Mapování} je technika, když paměť a registry externího zařízení se připojují ke sběrnici,
což umožňuje přístup k tomuto zařízení. Mohou se připojovat pouze jednotlivé adresy, nebo celé
rozsahy pamětí procesoru (obr. \ref{fig:memory_map_ex}).

\section{Časování a DMA}

Komunikace (konfigurace sběrnice a samotný přenos dat) na sběrnici vždycky zabírá určitý čas. Proto
opakovaný přístup ke sběrnici může zabrat cenný čas procesoru, který mohl by být stráven na jiné
užitečné úlohy.

Typickým řešením tomu je tzv. přímý přístup do pamětí (\textit{direct-memory access} --
\textit{DMA}). Je to speciální obvod (obvykle uvnitř procesoru), který slouží k asynchronnímu
přenosu dat. Princip DMA spočívá v tom, že procesor nastaví adresu, odkud DMA má číst data, a
cílenou adresu, kam DMA má zapisovat. Pak procesor spustí transakci, a DMA začne přenášet nezávislé
na procesoru. To umožňuje uvolnit CPU od přenosu dat, díky čemuž může plnit jiné úkoly. Přitom v
některých systémech DMA může dokonce přenášet na vyšších frekvencích než CPU.

V některých systémech ale existuje zjednodušený DMA, který sice je schopen přenášet data rychleji
než procesor, procesor bude pozastaven. Příkladem může být např. herní systém NES, kde přenos dat
mezi grafickou a zvukovou jednotkou pozastaví vykonávání instrukcí procesoru, dokud transakce nebude
dokončená.
